\documentclass[11pt,a4paper]{article}
\usepackage[T1]{fontenc}
\usepackage[utf8]{inputenc}
\usepackage[margin=2.5cm]{geometry}
\usepackage{enumitem}
\usepackage{array}

\title{Final Project Proposal\\[0.3em]
\large From Orbit to 3D: Single-View vs.\ Stereo Deep Learning\\for Satellite Terrain Reconstruction}
\author{Arthur Dohnert\\
\small INF01081 / CMP620 -- Deep Learning for Visual Computing, UFRGS, 2026/2}
\date{}

\begin{document}
\maketitle

\noindent\textbf{Group:} individual project (no partner).\\
\textbf{Option:} 1st and only option, a topic outside the suggested list.

\section*{Description}

Build a 3D model of a city (buildings, trees and streets) using only satellite images.
Two deep learning approaches will be compared on the same places:
\begin{itemize}
  \item \textbf{Path A:} estimate the height of every pixel from a \emph{single} image (U-Net with a pretrained encoder).
  \item \textbf{Path B:} estimate the height from \emph{two} images of the same place taken from different angles (stereo network, PSMNet).
\end{itemize}
The reconstructed terrain will be shown in an interactive 3D viewer, similar to the 3D mode of Google Earth, and used to plan the route of a rover, as is done for Mars rovers using orbital images.

\section*{Goals}

\begin{enumerate}
  \item Measure how much the second image improves the reconstruction (height error in meters).
  \item Find out in which cases stereo helps and in which a single image is enough.
  \item Build a 3D viewer of the reconstructed terrain, to be used as the live demo.
  \item Plan a rover route on the reconstructed terrain and check it against the real terrain.
\end{enumerate}

\section*{Resources}

\begin{itemize}
  \item \textbf{Data:} IEEE GRSS Data Fusion Contest 2019 (US3D), Track 2: about 4{,}300 stereo pairs of satellite images of Jacksonville and Omaha, with LiDAR height as ground truth.
  \item \textbf{Compute:} PCAD cluster (INF/UFRGS, RTX 4090 nodes) for training; laptop GPU (RTX 3050) for tests.
  \item \textbf{Tools:} Python, PyTorch, segmentation-models-pytorch, three.js (viewer).
\end{itemize}

\section*{Steps}

\begin{enumerate}
  \item Download and inspect the dataset; split train, validation and test by location. \emph{(done)}
  \item Run Path A and Path B end to end on a small part of the data, with a fair evaluation script. \emph{(done)}
  \item Train Path A and Path B on the full dataset on PCAD.
  \item Improve Path B with a pretrained stereo network (RAFT-Stereo) and add a classical stereo baseline (SGM).
  \item Compare A and B on the test set and analyze when each one works better.
  \item Finish the 3D viewer and the rover route planning.
  \item Prepare the slides and the live demo.
  \item Final presentation.
\end{enumerate}

\section*{Timetable}

\begin{center}
\renewcommand{\arraystretch}{1.3}
\begin{tabular}{|l|c|c|c|c|c|c|}
\hline
Step & Nov 3--9 & Nov 10--16 & Nov 17--23 & Nov 24--30 & Dec 1--6 & Dec 7 \\
\hline
1. Dataset          & x &   &   &   &   &   \\ \hline
2. First pipelines  & x &   &   &   &   &   \\ \hline
3. Full training    &   & x &   &   &   &   \\ \hline
4. Improve Path B   &   &   & x &   &   &   \\ \hline
5. Comparison       &   &   & x & x &   &   \\ \hline
6. Viewer and rover &   &   &   & x &   &   \\ \hline
7. Slides and demo  &   &   &   &   & x &   \\ \hline
8. Presentation     &   &   &   &   &   & x \\ \hline
\end{tabular}
\end{center}

\noindent Experiments end on November 30. The last week is only for the slides and the demo.

\end{document}
